\begin{textsample}{POS Dim 3 – go – Score 24.00 – go/go\_inf\_11.txt}  \label{ex:f3_pos_045}
3.1.
O Eu, o Outro e o Nós A \textbf{criança}, desde \textbf{bebê}, é capaz de explorar o espaço e tudo que contém nele e se \textbf{interessa} em interagir com as pessoas, \textbf{demonstrando} \textbf{afeto}, \textbf{desejo} e vontade.
Segundo Guimarães ( 2018 ), geralmente, as \textbf{crianças} são caracterizadas pelos \textbf{adultos} a partir daquilo que lhes faltam.
Não andam, não falam, estabelecem relações frágeis com os objetos, com outras \textbf{crianças} e \textbf{adultos}.
Assim, os professores, para romper com essa visão, precisam identificar : o que sabem as \textbf{crianças} desde \textbf{bebês}?
Como se comunicam?
Como expressam seus \textbf{desejos}?
O que sentem?
Quem são?
No intuito de conhecê -las e perceber suas potencialidades, preferências, \textbf{desejos} e interesses.
Nesse sentido, este campo dá visibilidade para as experiências vivenciadas pela \textbf{criança} no que diz respeito aos sentidos, saberes e conhecimentos produzidos sobre si e sobre o outro, numa perspectiva de interdependência e interrelação que possibilita tomar consciência da existência de um nós, de um grupo social no qual se vive um processo constante de humanização.
Essas experiências estão diretamente relacionadas a ações do cotidiano, como : a construção de vínculos, a afetividade, o autoconhecimento, a comunicação e o pertencimento a um determinado grupo social.
Portanto, esse campo traz à tona a importância da intencionalidade pedagógica no desenvolvimento dessas práticas no dia a dia das instituições de Educação \textbf{Infantil}.
É necessário compreender que o processo de humanização ocorre ao longo da vida, na medida em que o homem atua sobre a realidade, apropriando -se dos elementos da cultura e esse processo é particularmente intenso na infância.
A \textbf{criança} nasce com as condições físicas e biológicas para se desenvolver, mas é na relação estabelecida com outros sujeitos, em situações vividas coletivamente, que ela se apropria da experiência humana, desenvolve habilidades e qualidades como andar, falar, pensar, imaginar, que são próprias do ser humano.
Ter clareza que o processo de humanização não é natural, mas precisa ser aprendido e apropriado pela \textbf{criança}, faz total diferença para o trabalho da instituição de Educação \textbf{Infantil} e, nesse contexto, exige que seja planejado intencionalmente.
A BNCC ( BRASIL, 2017 ) apresenta este campo de experiências abordando os conceitos centrais para a compreensão dos processos de humanização e da formação das \textbf{crianças} em sua integralidade, como as interações, a autonomia, o autocuidado e a identidade, conforme o quadro a seguir : Quadro 10 – Campo de experiências O Eu o Outro e o Nós É na interação com os pares e com \textbf{adultos} que as \textbf{crianças} vão constituindo um modo próprio de agir, sentir e pensar e vão descobrindo que existem outros modos de vida, pessoas diferentes, com outros pontos de vista.
Conforme vivem suas primeiras experiências sociais ( na família, na instituição escolar, na coletividade ), constroem percepções e questionamentos sobre si e sobre os outros, diferenciando -se e, simultaneamente, identificando -se como seres individuais e sociais.
Ao mesmo tempo que participam de relações sociais e de \textbf{cuidados} pessoais, as \textbf{crianças} constroem sua autonomia e senso de autocuidado, de reciprocidade e de interdependência com o meio.
Por sua vez, na Educação \textbf{Infantil}, é preciso criar oportunidades para que as \textbf{crianças} entrem em contato com outros grupos sociais e culturais, outros modos de vida, diferentes atitudes, técnicas e rituais de \textbf{cuidados} pessoais e do grupo, costumes, celebrações e narrativas.
Nessas experiências, elas podem ampliar o modo de perceber a si mesmas e ao outro, valorizar sua identidade, respeitar os outros e reconhecer as diferenças que nos constituem como seres humanos [ grifos nossos ].
Fonte : Base Nacional Comum Curricular ( BRASIL, 2017 ).
Interações A constituição do eu é um processo tipicamente humano.
A \textbf{criança}, desde seu nascimento, está inserida em um grupo social caracterizado por uma herança cultural, que compreende conhecimentos, saberes, \textbf{hábitos}, costumes, objetos, formas de pensar, linguagens.
O processo de humanização vai se constituindo conforme a \textbf{criança} atua na realidade, se apropriando dos elementos dessa cultura que é partilhada pelas pessoas com as quais convive ( BARROS E VICENTINI, 2017 ).
Nesse contexto, compreende -se que o ser humano se constitui nas interações que estabelece com os sujeitos, com os artefatos, com a natureza e consigo mesmo, pois elabora sua existência nas vivências e experiências sociais.
Desde muito cedo, a \textbf{criança} é capaz de interagir com pessoas, explorar espaços e objetos, levantar hipóteses e elaborar explicações sobre fatos e situações que vivencia.
As interações são as vias de comunicação que conectam os sujeitos ao mundo, sendo um modo privilegiado de agir, de relacionar -se, de produzir sentidos e significados.
Dessa forma, as interações possibilitam a articulação e exploração das linguagens que a \textbf{criança} utiliza para se expressar, colocar -se na relação com o outro, aprender e desenvolver.
Para ocorrerem de forma integral e serem qualificadas, envolvem a afetividade, a linguagem, o pensamento, a corporeidade e a sociabilidade ( BEBER, 2018 ).
A \textbf{criança} aprende e desenvolve nas interações estabelecidas na companhia de idosos, \textbf{adultos}, jovens, adolescentes e \textbf{crianças}, no confronto de \textbf{gestos}, falas, ações, modificando sua forma de pensar, agir, sentir.
Nessas interações, cada \textbf{criança} apresenta um ritmo e uma forma própria de se envolver, de expressar suas emoções e \textbf{curiosidades} ( BRASIL, 2009 ).
A instituição de Educação \textbf{Infantil} é um dos espaços privilegiados para as interações com o outro, inclusive com seus pares de idade, pois é um grupo social mais amplo, diferente da família, onde os encontros ocorrem com regularidade, possibilitando a convivência, a troca de ideias, a construção de regras e as negociações.
Para qualificar essas interações, é fundamental o planejamento intencional de contextos significativos de aprendizagens, que permitam diálogos, questionamentos, cuidados de si e do outro, exploração de espaços naturais e construídos, de objetos e materiais disponibilizados às \textbf{crianças}.
Ampliar as oportunidades de experimentar, explorar, participar e vivenciar novas situações, novos contextos e aprendizagens, na convivência com outras \textbf{crianças} e \textbf{adultos}, é fundamental para que sentidos, saberes e conhecimentos sejam compartilhados, possibilitando que conheçam aspectos culturais dos grupos sociais com os quais convivem e aprendam a identificar suas preferências e características, cuidando de si e do outro, numa dimensão afetiva, \textbf{lúdica} e de diálogo.
Autonomia As formas de ser e se relacionar com o mundo são caracterizadas pelos costumes presentes na cultura de um grupo social.
A partir da comunicação por meio de diferentes linguagens, da observação, da participação em variadas situações, da exploração do espaço, a \textbf{criança} se apropria de comportamentos e \textbf{hábitos} tipicamente humanos.
Tais comportamentos e \textbf{hábitos} são marcados por uma cultura, uma vez que, o jeito de dormir, de se alimentar, de cumprimentar o outro, de participar de uma festividade, de como se portar em diferentes lugares, varia a depender do contexto histórico em que essas ações são realizadas.
Assim, conforme a \textbf{criança} compreende e se apropria desses \textbf{hábitos}, começa a desenvolver e a assumir uma autonomia pessoal e corporal.
É incontestável a importância da aprendizagem dos \textbf{cuidados} pessoais.
Ter consciência e controle do corpo, aprender a se vestir, a colocar seus sapatos, ter iniciativa de tomar água, de se alimentar sozinho, são ações que à primeira vista parecem simples, mas para a \textbf{criança} se constituem em desafios e grandes conquistas.
Aprender a se relacionar, respeitando regras de convivência social, também é outro aspecto importante da autonomia.
A experiência da vida coletiva ensina as \textbf{crianças} a colaborarem entre si e a compreenderem os valores que orientam as atitudes de colaboração, respeito ao outro, dignidade, integridade, liberdade, igualdade, inviolabilidade da vida humana, que são a base para a construção das relações da vida em sociedade ( BRASIL, 2009 ).
Esses valores, que são estabelecidos em determinado grupo social, sustentam as relações dos sujeitos na comunidade.
Os valores são as referências para se compreender o que é culturalmente aceito ou não na conduta dos sujeitos nas interações sociais.
A instituição educacional é um espaço privilegiado para essas aprendizagens ao permitir a negociação e o compartilhamento de diferentes modos de ser e de agir.
Ao comunicar com os sujeitos suas experiências, suas ideias, seus pensamentos, seus \textbf{desejos}, suas necessidades, seus sentimentos, a \textbf{criança} se constitui e vai constituindo outro grupo social para conviver.
Aprender a cuidar de si, do outro, das relações, são ações que precisam ser ensinadas.
A \textbf{criança} aprende na relação com os \textbf{adultos} e com as \textbf{crianças} mais experientes.
Essas aprendizagens precisam ser mediadas em situações de diálogo, em que ela possa expressar sentimentos, sensações e não com atitudes impositivas, controladoras e autoritárias.
Dessa forma, a \textbf{criança} desenvolve sua autonomia, fazendo, experimentando, errando, \textbf{manipulando}, observando e \textbf{imitando}, nas situações cotidianas vivenciadas em casa, na instituição educacional, enfim, em diferentes espaços.
Autocuidado O Parecer das DCNEB ( BRASIL, 2010 ) estabelece que o trabalho de uma instituição educacional tem sua centralidade na formação da pessoa na sua essência humana.
Nesta perspectiva, é fundamental considerar as dimensões do educar e cuidar de forma inseparável, sendo necessário compreender que enfrentar o desafio de educar, neste mundo complexo, exige relações pautadas no acolhimento de todos, com respeito e com atenção.
Assim, > Em cada \textbf{criança} [... ] há uma criatura humana em formação e, nesse sentido, cuidar e educar são, ao mesmo tempo, princípios e atos que orientam e dão sentido aos processos de ensino, de aprendizagem e de construção da pessoa humana em suas múltiplas dimensões ( BRASIL, 2010, p. 18 ).
Na instituição de Educação \textbf{Infantil}, tratar dessas ações de \textbf{cuidado} exige uma disponibilidade à \textbf{escuta}, como algo que me conecta ao outro, não só pela audição, mas também pela visão, pelo tato, pela observação, pela sensibilidade.
Quando alguém se propõe a escutar o outro, está presente uma \textbf{curiosidade}, um \textbf{desejo}, uma dúvida, um interesse.
Exige que certezas, julgamentos sejam suspensos, para legitimar e valorizar os pensamentos, as palavras, as expressões e as necessidades da \textbf{criança} ( RINALDI, 2014 ).
Destaca -se aqui a experiência do olhar.
O \textbf{bebê} busca com o olhar no outro a sensação de confiança para dar seguimento a uma ação de exploração ao tocar um objeto, ao se movimentar em direção a algo que lhe chame a atenção.
O papel do professor é \textbf{acolher} esse olhar, perceber e significar essa interação que envolve necessidades, sentimentos e emoções.
Para tanto, é necessário que o professor esteja atento às suas próprias ações, se preocupando com a qualidade ética do \textbf{cuidado}, refletindo sobre o sentido do seu modo de fazer, sentir, olhar as \textbf{crianças}, podendo assim, sustentar as experiências delas consigo e com o mundo ( GUIMARÃES, 2011 ).
O \textbf{cuidado}, nessa dimensão ética, envolve as ações de proteção, relativas às necessidades físicas e biológicas da \textbf{criança}, como o sono, a alimentação, a \textbf{higiene} pessoal.
Nesse contexto, também é importante o desenvolvimento das potencialidades da \textbf{criança}, a construção da \textbf{autoconfiança} e da autonomia, na qual são valorizadas as iniciativas próprias, a descoberta das possibilidades e limites de seu corpo, de suas ações, construindo capacidades de escolher, comunicar, tendo oportunidades de se conhecer ( GUIMARÃES, 2011 ).
Essa dimensão do \textbf{cuidado} também deve se fazer presente nas situações de conflito entre as \textbf{crianças}, quando ocorre manifestações agressivas, de desrespeito ao outro.
É importante que se compreenda o contexto, a ação que estava sendo desenvolvida no momento, a organização do tempo e do espaço, as relações e não apenas a ação isolada.
Permitir que as \textbf{crianças} manifestem o que e como se sentem, pensam, não significa reforçar essas atitudes, e sim estabelecer um diálogo que promova uma reflexão sobre os motivos e os resultados da ação, de uma forma em que se sinta acolhida, protegida e não rejeitada por castigos e isolamentos.
As relações de educação e \textbf{cuidado} envolvem convivência e respeito pelo outro, em seus saberes, conhecimentos, \textbf{hábitos}, costumes, crenças, ritmos pessoais, características físicas etc., com perspectiva de superar ações automatizadas e mecânicas, em que as \textbf{crianças} são tratadas como objetos nos atos de trocar \textbf{fraldas}, dar banho, alimentar, mediar situações conflituosas ou de desenvolver outras atividades.
Ao contrário, se há diálogo, olho no olho, toque com \textbf{afeto}, elas desenvolvem a \textbf{autoconfiança}, o que permite dar conta de si mesma, assim como expressar -se de forma segura e prazerosa ( GUIMARÃES, 2011 ).
Identidade Identidade é a consciência que uma pessoa tem dela própria, tornando -a alguém diferente das outras.
É constituída a partir de todas as experiências vivenciadas nas interações com o outro, com os objetos e os espaços da cultura.
As características físicas demarcam algumas diferenças, como a cor dos olhos, dos cabelos, da pele, mas é nas relações concretas com um grupo social e sua cultura que a \textbf{criança} passa a existir.
Receber um nome, ser filho(a) de pessoas específicas, nascer em um país, cidade, bairro onde vivencia um contexto social e cultural, a torna uma \textbf{criança} diferente das outras ( MELLO, 2010 ).
O processo de constituição da identidade ocorre desde antes do nascimento.
As experiências \textbf{sensoriais} e motoras, marcadas pela afetividade, iniciam esse processo.
Por \textbf{volta} de um ano e meio, em frente a um espelho, a \textbf{criança} já é capaz de perceber sua própria imagem, distinguindo -se de outras.
Com o desenvolvimento da linguagem oral, esse processo se amplia, associando -se ao desenvolvimento da autopercepção.
Primeiro ela se refere a si mesma pelo seu nome e só depois começa a fazer uso de pronomes - eu, mim.
É na convivência com as pessoas que ela constrói vínculos afetivos e passa por um processo de reconhecimento do outro, de constatação das diferenças entre os sujeitos e daquilo que a identifica e a diferencia.
Nesse contexto, é de suma importância a compreensão da alteridade, que expressa a qualidade do que é o “ outro ”, do que é o diferente.
A consciência do “ eu ” só é possível por meio do reconhecimento do “ outro ”, o que implica que um sujeito seja capaz de se colocar no lugar do outro, numa interação fundamentada no diálogo e na valorização das diferenças.
Na instituição de Educação \textbf{Infantil}, para pensar nesse “ outro ” \textbf{criança}, é necessário reconhecer as múltiplas formas de ser \textbf{criança} e de viver a infância, frente às diferentes realidades sociais que se apresentam.
Assim, considerar a alteridade da infância “ requer o resgate com a experiência humana ” ( OLIVEIRA, 2004, p.09 ).
Para que este resgate aconteça, é necessário que se compreenda as múltiplas linguagens que a \textbf{criança} usa para se expressar, que se aprenda a escutar, registrar, representar suas vozes e seus movimentos.
Destaca -se o quanto é fundamental a constituição do grupo no espaço institucional, pois nele a \textbf{criança} tem a possibilidade de exercitar e confrontar as diferenças, aprendendo a difícil tarefa de conviver com divergências e conflitos.
O primeiro grupo social da \textbf{criança} é a família, compreendida em suas diferentes configurações para além da composta por pai, mãe e filho, na qual ocorre a construção de vínculos afetivos duradouros que são ampliados nas instituições educativas.
Essas relações, estabelecidas nos diferentes espaços sociais no início da vida, são muito importantes para o desenvolvimento da identidade, pois pertencer a um grupo social é fazer parte de um coletivo, pertencer a um lugar, a uma comunidade.
É conviver diariamente em um espaço onde se vivencia significados, tradições, ações e formas de pensar.
É compartilhar um território, uma história, uma linguagem, uma forma de ser e sentir.
Enquanto pertencente a um grupo, a \textbf{criança} se diferencia como um ser individual que tem características físicas e de personalidade, vontades e necessidades próprias.
A partir das interações, ela percebe que no espaço institucional é constituído um grupo, que é diferente de outros.
Essas relações coletivas são fundamentais para a \textbf{criança} tomar consciência da existência de um “ nós ”, no qual se estabelecem relações que afetam ao grupo e a ela mesma.
As identidades culturais que compõem os sujeitos das instituições precisam ser valorizadas e afirmadas, admitindo -se suas diversidades, compreendidas nos costumes, nas tradições, na culinária, na religião, na linguagem, na organização familiar.
Reconhecê -las está para além da aceitação, segundo Barbosa ( 2009 ), implica em compreender que as diferenças são o que são, e não vão deixar de ser, e sim, serão colocadas em relação, em negociação.
Assim, considera -se que > Somos todos diversos em nossa biologia, > em nossa experiência, nosso modo de ser, > em nossa cultura.
A diversidade biológica e a diversidade cultural são a regra da > espécie humana e também no planeta [... ] > pensar em uma educação em que assumamos que a diferença também está em nós, > ou ainda, que a diferença somos nós e que > a escola é um importante lugar de encontro/confronto com as diferenças e a construção de suas negociações.
Os sujeitos, as > culturas, as religiões, as etnias já existiam > antes da escola, cabe à escola incluí -los, > para que ela possa permanecer existindo. > ( BARBOSA, 2009, p.13 ) Os elementos culturais que constituem a diversidade das \textbf{crianças} e de suas famílias, relacionados à etnia, ao gênero, à classe social, ao espaço geográfico, precisam ser explorados e valorizados no planejamento de contextos significativos de aprendizagens, de forma que esses elementos sejam visíveis nos espaços, no cardápio, na organização do momento do sono, em projetos investigativos etc., permitindo que se reconheçam e construam uma imagem positiva de si e de seus grupos de pertencimento.
As DCNEI ( BRASIL, 2009 ) preconizam que as práticas pedagógicas das instituições de Educação \textbf{Infantil} promovam : uma visão plural do mundo pela \textbf{criança} ; a ampliação do seu olhar para diferentes povos e culturas, uma relação positiva e apropriação das contribuições histórico-culturais dos povos indígenas, afrodescendentes, asiáticos, europeus e de outros países da América.
A instituição também precisa se preocupar com as condições de acesso, efetiva participação e aprendizagens a todas as \textbf{crianças}, com base na igualdade de oportunidades, não permitindo que sejam excluídas em razão de alguma deficiência física ou cognitiva, considerando que a educação é um direito inalienável, previsto na Constituição Federal ( BRASIL, 1988 ), e deve ser assegurada sem nenhum tipo de distinção.
A Proposta Pedagógica precisa cumprir sua função social, valorizando as diferenças e atendendo as necessidades educacionais específicas das \textbf{crianças}.
Conforme consta nas DCNEI ( BRASIL, 2009 ), no planejamento pedagógico é importante pensar em estratégias que permitam às \textbf{crianças} com necessidades educativas especiais vivenciar o contexto da instituição em sua amplitude, pensando os espaços, os materiais, os \textbf{brinquedos}, as formas de comunicação, os procedimentos que viabilizem às \textbf{crianças} serem ativas, possibilitando a sua ação nas \textbf{brincadeiras} e interações com as outras \textbf{crianças}.
Assim, esse campo de experiência propõe que a instituição de Educação \textbf{Infantil} seja um espaço onde as \textbf{crianças} possam constituir suas identidades, imersas em interações qualificadas, vivenciadas em contextos diferentes de aprendizagens que permitam o estabelecimento de relações com pessoas, objetos e a apropriação de conhecimentos produzidos pela humanidade, ampliando suas possibilidades de compreensão e ação no mundo.

% matched lemmas: acolher, adulto, afeto, autoconfiança, bebê, brincadeira, brinquedo, criança, cuidado, curiosidade, demonstrar, desejo, escuta, fralda, gesto, higiene, hábito, imitar, infantil, interessar, lúdico, manipular, sensorial, volta
\end{textsample}
